\documentclass[12pt, a4paper, twoside]{article}
\usepackage[utf8]{inputenc}
\usepackage[T1]{fontenc}
\usepackage[brazil]{babel}
\usepackage[top=3.0cm, left=3.0cm, bottom=2.0cm, right=2.0cm]{geometry}
\usepackage{setspace}
\usepackage{graphicx}
\usepackage{booktabs}
\usepackage{tabularx}
\usepackage{multirow}
\usepackage{float}
\usepackage{amsmath}
\usepackage{amssymb}
\usepackage[alf]{abntex2cite}

\graphicspath{{imagens/}}

\newcommand{\mesatual}{%
  \ifcase\month\or
    Janeiro\or Fevereiro\or Março\or Abril\or Maio\or Junho\or
    Julho\or Agosto\or Setembro\or Outubro\or Novembro\or Dezembro
  \fi
}

\begin{document}
\onehalfspacing

% CAPA
\begin{titlepage}
    \begin{center}
        \includegraphics[height=2cm]{logo-ppgia.jpeg}
        \vspace{0.5cm}

        % \includegraphics[height=2.8cm]{logo-ufrpe.png}
        % \hspace{2cm}
        % % \includegraphics[height=2cm]{logo-ppgia.jpeg}
        % % \vspace{0.5cm} \\
        \textbf{UNIVERSIDADE FEDERAL RURAL DE PERNAMBUCO} \\
        \textbf{PRÓ-REITORIA DE PESQUISA E PÓS-GRADUAÇÃO} \\
        \textbf{PPG EM INFORMÁTICA APLICADA} \\
        \vspace{5cm}
        \textbf{\Large PREVISÃO DE SÉRIES TEMPORAIS FINANCEIRAS: UMA AVALIAÇÃO DE ESTRATÉGIAS DE OTIMIZAÇÃO PARA MODELOS BASEADOS EM LSTM E GRU} \\
        \vfill
        \begin{flushright}
            \textbf{ALUNO:} Arthur Lima Silva \\
            \textbf{ORIENTADOR:} George Gomes Cabral
        \end{flushright}
        \vfill
        \textbf{RECIFE-PE -- \mesatual/\the\year}
    \end{center}
\end{titlepage}

% IDENTIFICAÇÃO DO PLANO
\newpage
\begin{center}
\textbf{Identificação do plano}    
\end{center}

\vspace{0.5cm}

\noindent \textbf{a) Nome do discente:} Arthur Lima Silva \\
\textbf{b) Bolsista Sim ( ) Não (X) Agência:} - \\
\textbf{c) Nível (mestrado/doutorado):} mestrado \\
\textbf{d) Entrada (ano/semestre):} 2026.1 \\
\textbf{e) Orientador:} George Gomes Cabral \\
\textbf{f) Co-orientadores:} - \\
\textbf{g) Projeto de pesquisa do(a) orientador(a) ao qual está vinculado o plano (projeto guarda-chuva):} Previsão de Séries Temporais para o Mercado de Ações: Integração de Abordagens Estatísticas, Aprendizado de Máquina e Medidas Probabilísticas de Risco \\
\textbf{h) Área de concentração:} Inteligência Artificial \\
\textbf{i) Linha de pesquisa:} Aprendizado de Máquina e Previsão de Séries Temporais

\newpage

% 1 INTRODUÇÃO
\section{Introdução}

A previsão do comportamento dos mercados financeiros é um desafio bastante estudado pelas áreas de finanças quantitativas e inteligência artificial. Historicamente, a motivação primária para a modelagem de séries temporais financeiras centrou-se na maximização de lucros e na gestão de portfólios por parte de investidores e instituições privadas. Contudo, a relevância dessa área de estudo transcende o ganho de capital isolado. Como destacado por \citeonline{li_accurate_2024}, o mercado de ações não apenas reflete o crescimento econômico de um país, mas também serve como base fundamental para a formulação de políticas macroeconômicas.

Nesse contexto, o mercado de ações atua como um termômetro do ânimo dos investidores e do comportamento humano, visto que flutuações no mercado de ações frequentemente antecedem ou acompanham variações no Produto Interno Bruto (PIB) das nações, conforme argumentado por \citeonline{ramkumar_development_2023}. Essa visão é corroborada por \citeonline{sun_novel_2024}, que posicionam o mercado de ações como um ponto focal imprescindível para que os Estados regulem tendências econômicas e mitiguem o impacto de crises sistêmicas.

Além de seu papel como indicador macroeconômico, a previsibilidade financeira afeta diretamente infraestruturas críticas e a tomada de decisão em setores estratégicos. O estudo de \citeonline{hosseini_pattern-based_2025} exemplifica essa dinâmica ao demonstrar que, em mercados como o de energia e eletricidade na Europa, o sucesso operacional e o equilíbrio entre oferta e demanda são fortemente influenciados pela capacidade preditiva dos agentes envolvidos. Portanto, o aprimoramento de modelos de previsão financeira justifica-se não apenas pela mitigação de riscos e contenção de perdas em portfólios de investimento, mas como uma ferramenta essencial de governança, estabilidade econômica e suporte à tomada de decisão.

A complexidade inerente às séries temporais financeiras - caracterizadas por alta volatilidade, não estacionariedade e grande presença de ruído - exige a evolução contínua das metodologias de previsão. Historicamente, abordagens estatísticas tradicionais e algoritmos clássicos de \textit{Machine Learning} estabeleceram as fundações dessa área. Trabalhos recentes \cite{mahjouby_predicting_2024, sharma_stock_2023} ainda documentam o uso de modelos como ARIMA, SVM, \textit{Random Forest} e \textit{Ensembles} de árvores de decisão. Embora esses trabalhos estabeleçam \textit{baselines} sólidas, e o XGBoost ainda apresente resultados competitivos entre as abordagens clássicas, a literatura aponta que esses métodos frequentemente esbarram em limitações quando expostos a padrões ocultos profundos e relações não-lineares severas.

Para superar essas barreiras, o estado da arte tem convergido para o uso de \textit{Deep Learning}, com forte predominância das redes \textit{Long Short-Term Memory} (LSTM), devido à sua capacidade de modelar dependências temporais de longo prazo. Os trabalhos de \citeonline{k_leveraging_2023} e \citeonline{miot_empirical_2020} são alguns dos que marcam essa transição ao evidenciar a superioridade de LSTM’s frente a RNN’s tradicionais e modelos ARIMA. Contudo, a pesquisa contemporânea reconhece que redes neurais individuais também possuem limitações, impulsionando o estudo de modelos híbridos.

Trabalhos recentes \cite{kiruba_risk_2025, ramkumar_development_2023} propõem a integração de LSTM com ARIMA, frequentemente aliados a estratégias de segmentação de dados para ganho de precisão. Na mesma linha de otimização através da combinação de técnicas, \citeonline{liu_application_2024} e \citeonline{choi_hybrid_2023} demonstraram que modelos híbridos são consistentemente superiores na seleção de ativos e mitigação de erros.

Atualmente, o aprimoramento dessas arquiteturas concentra-se em três frentes principais: o desenvolvimento de técnicas de pré-processamento e extração de características para reduzir ruídos e destacar padrões relevantes das séries temporais \cite{sun_novel_2024, li_accurate_2024, hosseini_pattern-based_2025, wan_tcdformer_2024}; a incorporação de mecanismos de atenção e arquiteturas capazes de modelar relacionamentos temporais e estruturais mais complexos, como \textit{Transformers}, \textit{Temporal Convolutional Networks} e \textit{Graph Attention Networks} \cite{biswas_dual_2025, luo_short-term_2024, hsu_fingat_2023, anghinoni_time_2019}; e a integração de informações multimodais, combinando dados históricos com notícias, indicadores econômicos e outras fontes externas para capturar fatores exógenos que influenciam o comportamento do mercado \cite{choi_hybrid_2023, wang_stock_2023, haryono_transformer-gated_2023, anbaee_farimani_adaptive_2024}.

Por fim, é fundamental pontuar o desafio metodológico sistêmico delineado por \citeonline{hu_fintsb_2025}. O artigo diagnostica que a comparação de abordagens entre trabalhos constantemente sofre com a falta de padronização (\textit{Standardization Deficit}), variações nos regimes de dados testados (\textit{Diversity Gap}) e a omissão de fricções do mundo real (\textit{Real-World Mismatch}).

A partir do mapeamento do estado da arte, evidencia-se que a superação das barreiras de predição financeira não reside exclusivamente na criação de redes neurais mais profundas, mas na capacidade de fundir informações contextuais (multimodalidade), técnicas de remoção de ruído (decomposição) e quantificação de incertezas. Sendo assim, o presente projeto de dissertação propõe responder ao seguinte problema de pesquisa: \textbf{Como melhorar o desempenho e a robustez de modelos baseados em LSTM através da extração avançada de características e da integração de dados externos multimodais para a classificação e previsão de séries temporais no mercado de ações?}

\newpage

\section{Objetivos}

\subsection{Geral}

Investigar, desenvolver e validar uma arquitetura híbrida para a previsão de séries temporais do mercado de ações, baseada em modelos recorrentes (como LSTM e GRU), avaliando e integrando diferentes estratégias de otimização metodológica, com o intuito de aumentar a precisão, a capacidade de generalização e a prevenção de riscos sob diferentes regimes de mercado.

\subsection{Específicos}

\begin{itemize}
    \item Aprofundar a análise do estado da arte referente à predição financeira e a modelos híbridos de \textit{Deep Learning}, com o intuito de identificar lacunas e selecionar estratégias de otimização mais promissoras a serem avaliadas na pesquisa;

    \item Estruturar um \textit{pipeline} de coleta e pré-processamento de dados de ativos financeiros (utilizando plataformas como \textit{Yahoo Finance} e a base FinTSB).

    \item Implementar e validar arquiteturas baseadas em redes recorrentes (LSTM e GRU), estabelecendo um \textit{baseline} de desempenho interno confiável para mensurar o impacto das otimizações subsequentes;

    \item Investigar, de forma comparativa, os efeitos da integração de diferentes estratétigas de aprimoramento do modelo base - avaliando vertentes da literatura como decomposição de seinais, extração de padrões, mecanismos de atenção ou fusão de variáveis exógenos -, a fim de determinar quais abordagens lidam melhor com o ruído e a complexidade das séries temporais financeiras;

    \item Avaliar o desempenho e a robustez das arquiteturas propostas, utilizando análises quantitativas e métricas de erro (MSE, RMSE, MAPE), submetendo os modelos a testes de estresse em cenários variados de tendência e volatilidade disponibilizados pelo \textit{benchmark} FinTSB.
\end{itemize}

\newpage

\section{Metodologia}

Esta pesquisa caracteriza-se por uma abordagem de natureza aplicada, quantitativa e experimental, direcionada ao desenvolvimento e validação de arquiteturas preditivas para o mercado de ações. A metodológia a ser seguida está estruturada em três fases sequenciais: (1) a definição do problema somada à revisão da literatura; (2) a implementação computacional, geração de resultados empíricos e análise de desempenho — que representa o núcleo prático deste estudo; e (3) a consolidação dos achados por meio da escrita científica. Na segunda fase, o foco experimental residirá na proposição de um módulo preditivo de base híbrida, construído a partir da integração das arquiteturas \textit{Long Short-Term Memory} (LSTM) e \textit{Gated Recurrent Unit} (GRU). A escolha por essas tecnologias não é arbitrária, fundamentando-se no atual estado da arte, que demonstra a eficácia superior de modelos recorrentes no processamento de sequências financeiras \cite{k_leveraging_2023, wang_stock_2023, haryono_transformer-gated_2023}.

Para viabilizar o treinamento e a avaliação inicial dos modelos, o processo de coleta de dados apoia-se em duas fontes de dados. A fonte primária de extração de séries temporais será a plataforma Yahoo Finance, que fornecerá os dados brutos diários (como abertura, fechamento, máximas, mínimas e volume) de ativos específicos selecionados para compor o cenário base de implementação. Contudo, para garantir que os resultados não sejam enviesados pela escolha de um grupo restrito de ações, a pesquisa utilizará um \textit{baseline} complementar. Uma forte candidata para essa função é a base de dados FinTSB \cite{hu_fintsb_2025}. Bases padronizadas de \textit{benchmark} como esta são metodologicamente valiosas pois categorizam e disponibilizam recortes de mercado sob diferentes regimes temporais — tais como tendências de alta (\textit{uptrends}), de baixa (\textit{downtrends}) e eventos de extrema volatilidade (\textit{black swan events}) —, oferecendo um acervo estruturado para testes rigorosos de generalização.

O aprimoramento do modelo preditivo, para além da arquitetura híbrida básica, envolverá a exploração de técnicas em consonância com as tendências recentes da literatura. Uma das possibilidades a serem investigadas durante o desenvolvimento é a otimização do pré-processamento. A literatura contemporânea sugere que a aplicação de técnicas de decomposição de sinais (como CEEMDAN e VMD) ajuda a isolar ruídos de tendências latentes \cite{sun_novel_2024, li_accurate_2024}. Em paralelo, pode-se avaliar a extração de características baseadas em padrões matemáticos de preços, conforme abordado por \citeonline{hosseini_pattern-based_2025}. Outra frente exploratória diz respeito à integração de Mecanismos de Atenção para forçar a rede a priorizar instâncias temporais críticas ou mapear a influência relacional entre diferentes ativos, inspirando-se em implementações em Redes Convolucionais Temporais \cite{biswas_dual_2025} e Redes de Atenção em Grafos \cite{hsu_fingat_2023}.

Ademais, outra via promissora de inovação metodológica a ser considerada é a exploração da multimodalidade. Sabendo que a precificação de ativos sofre forte impacto de variáveis exógenas, a pesquisa poderá investigar a fusão dos dados temporais com informações contextuais, como o sentimento extraído de notícias ou indicadores macroeconômicos. A viabilidade dessa integração tem sido amplamente demonstrada, seja por meio de arranjos baseados em \textit{Multilayer Perceptrons} \cite{choi_hybrid_2023}, seja pela incorporação de \textit{Encoders} da arquitetura \textit{Transformer} em conjunto com redes LSTM/GRU \cite{wang_stock_2023, haryono_transformer-gated_2023}.

Por fim, a avaliação experimental será conduzida por meio de análises comparativas quantitativas. O desempenho do modelo será mensurado utilizando métricas de erro consagradas no escopo de séries temporais financeiras, tais como o Erro Quadrático Médio (MSE), a Raiz do Erro Quadrático Médio (RMSE) e o Erro Absoluto Médio Percentual (MAPE), além de métricas de acurácia direcional que avaliam a classificação correta da tendência. É nesta etapa de validação que o uso de um \textit{baseline} complementar (como o FinTSB) desempenhará seu papel fundamental: ao submeter o modelo proposto a cenários padronizados e comparar seus resultados com outros modelos da literatura, será possível garantir estatisticamente que as melhorias obtidas representam um avanço estrutural genuíno e não um mero sobreajuste (\textit{overfitting}) aos dados inicialmente extraídos do Yahoo Finance.

\newpage

% 4. ORÇAMENTO
\section{Orçamento}

Não se aplica.

% 5 RESULTADOS ESPERADOS
\section{Resultados Esperados}

\begin{table}[H]
\begin{tabular}{lc}
\hline
\textbf{Categoria} & \textbf{Número esperado} \\ \hline
Artigo científico & 2 \\
Participação em congresso & 1 \\
Patente ou equivalente & 0 \\
Livro ou capítulo de livro & 0 \\
Dissertação de Mestrado & 1 \\ 
Tese de Doutorado & 0 \\ 
Participação em Editais & 0 \\ 
Outras (especificar) & 0 \\ \hline
\end{tabular}
\end{table}

% 6 CRONOGRAMA
\section{Cronograma do discente}

\begin{table}[H]
\centering
\resizebox{\textwidth}{!}{
\begin{tabular}{|l|c|c|c|c|c|c|c|c|}
\hline
\multirow{2}{*}{\textbf{ATIVIDADES}} & \multicolumn{8}{c|}{\textbf{SEMESTRES}} \\ \cline{2-9}
 & \textbf{1º} & \textbf{2º} & \textbf{3º} & \textbf{4º} & \textbf{5º} & \textbf{6º} & \textbf{7º} & \textbf{8º} \\ \hline
Créditos de disciplinas & X & X & X & & & & & \\ \hline
Créditos de orientação & & & & & & & & \\ \hline
Pesquisa Bibliográfica & X & X & & & & & & \\ \hline
Obtenção de dados & & X & X & & & & & \\ \hline
Análise de dados & & X & X & X & & & & \\ \hline
Proficiências (inglês/espanhol) & & X & & & & & & \\ \hline
Exame de qualificação* & & & & & & & & \\ \hline
Revisão e redação da tese/dissertação & & & X & X & & & & \\ \hline
Defesa da tese/dissertação & & & & X & & & & \\ \hline
\end{tabular}
}
\end{table}

\noindent * Só para doutorado.

% REFERÊNCIAS
\newpage
\bibliography{referencias}

\end{document}